% !TeX root = ../main.tex
\section{Deformation, distention, and pressure coupling}
\label{sec:finite-deformation-biot}

A saturated porous solid changes volume through deformation of its constituent
material and through changes in pore volume.  These mechanisms enter the stress
and the phase balances through different quantities.  We first identify those
quantities and then obtain the Biot coefficient from the mineral-volume
response at fixed pore pressure.

\subsection{Solid motion and phase mass}
\label{sec:solved-solid-state}

Let the solid motion map a reference position \(\mbf X\) to
\(\mbf x=\bs\chi(\mbf X,t)\).  Its deformation gradient and solid-reference
Jacobian are
\begin{equation}
 \mbf F=\pder{\bs\chi}{\mbf X} = \nabla_{\mbf X} \mbf{u}_s + \mbf{I},
 \qquad J=\det\mbf F>0,
 \label{eq:skeleton-kinematics}
\end{equation}
where $\mbf{u}_s$ is the solid displacement field. Following Drumheller \citep[Sec.~3]{drumheller2000}, we separate the true
motion of the solid material from the distention gradient,
\begin{equation}
 \mbf F=\mbf A\bar{\mbf F},
 \qquad \bar J=\det\bar{\mbf F},
 \qquad a=\det\mbf A,
 \qquad J=a\bar J.
 \label{eq:inelastic-multiplicative-split}
\end{equation}
Here \(\bar{\mbf F}\) is the true-deformation gradient of the solid
material, \(\bar J\) is its Jacobian, \(\mbf A\) is the distention
gradient, and \(a\) is the distention.  We choose a reference state with
\(J_0=\bar J_0=a_0=1\).  The ratio \(\bar J\) measures the change in
mineral volume, and the distention measures the additional volume change of
the porous structure: increasing \(a\) at fixed mineral volume increases
pore volume.

Let \(\bar\rho_s\) be the solid intrinsic density and \(\phi_s\) its volume
fraction.  The partial density per unit current mixture volume is
\begin{equation}
 \rho_s=\phi_s\bar\rho_s.
 \label{eq:solid-volume-fraction-relation}
\end{equation}
For a solid moving with velocity \(\mbf v_s\), mass conservation
in spatial and solid-reference forms is
\begin{align}
 \left.\pder{\rho_s}{t}\right|_{\mbf x}
 +\nabla_{\mbf x}\mathbin{\cdot}(\rho_s\mbf v_s)&=0,
 \label{eq:spatial-solid-mass-balance}\\
 \left.\pder{(J\rho_s)}{t}\right|_{\mbf X}&=0.
 \label{eq:reference-solid-mass-balance}
\end{align}
For the uniform reference state used here,
\(J\rho_s=\rho_{s0}=\phi_{s0}\bar\rho_{s0}\).  The true motion preserves
mineral mass, giving
\begin{equation}
 \bar J\equiv\frac{\bar\rho_{s0}}{\bar\rho_s},
 \qquad \phi_s=\frac{\phi_{s0}\bar J}{J},
 \qquad \phi_s a=\phi_{s0}.
 \label{eq:distension}
\end{equation}
Thus \(\bar\rho_s\) determines the mineral volume ratio
\(\bar J=\bar\rho_{s0}/\bar\rho_s\), whereas
\(\rho_s/\bar\rho_s=\phi_s\) is the fraction of the current mixture volume
occupied by solid.  For a fully saturated mixture, the fluid volume fraction
is therefore \(\phi_f=1-\phi_s\).

\subsection{Pressure and the two effective stresses}
\label{sec:intrinsic-pressure}

Let \(\psi_s\) denote the solid Helmholtz free energy per unit mass, with skeleton
deformation, intrinsic density, and constitutive internal state as its
arguments.  Local equilibrium sets the solid thermodynamic pressure,
\(\bar\rho_s^2\partial\psi_s/\partial\bar\rho_s\), equal to the pore pressure
\citep[eq.~(27c)]{fosterxu2025},
\begin{equation}
 p=\bar p_f=\bar p_s
 =\bar\rho_s^2\left.\pder{\psi_s}{\bar\rho_s}\right|_{\mbf F},
 \label{eq:local-pressure-equilibrium}
\end{equation}
where \(\bar p_f\) and \(\bar p_s\) are the intrinsic fluid and solid
thermodynamic pressures, respectively, and any internal variables are held
fixed.  The derivative in \eqref{eq:local-pressure-equilibrium} is evaluated
at fixed \(\mbf F\), so it equates the scalar pressure conjugate to a change
in intrinsic density without varying the skeleton deformation.  The separate
deformation derivative of the energy may therefore include a deviatoric stress
response.

Write \(W_s=\rho_{s0}\psi_s\) per unit reference mixture volume.  In terms
of \(\bar J\), the pressure relation \eqref{eq:local-pressure-equilibrium}
becomes
\begin{equation}
 p=-\frac{1}{\phi_{s0}}
 \left.\pder{W_s}{\bar J}\right|_{\mbf F}.
 \label{eq:intrinsic-pressure-free-energy}
\end{equation}
The independent arguments of \(W_s\) are \(\mbf F\) and
\(\bar\rho_s\), equivalently \(\mbf F\) and \(\bar J\), with internal
variables held fixed.  Solving this pressure relation determines
\(\bar J(\mbf F,p)\).  Substituting this solution into
\(W_s(\mbf F,\bar J)\) defines the fixed-pressure potential
\(W^{\prime\prime}(\mbf F,p)
\equiv W_s(\mbf F,\bar J(\mbf F,p))\).  The pressure Legendre transform, a
change of independent variable from intrinsic density to pressure, defines
the potential
\begin{equation}
 W^{\prime}(\mbf F,p)
 =W^{\prime\prime}(\mbf F,p)+\phi_{s0}p\bar J(\mbf F,p).
 \label{eq:pressure-legendre-potential}
\end{equation}
The single-prime first Piola stress is the deformation derivative of
\(W_s\) at fixed intrinsic density, equivalently the derivative of the
transformed potential \(W^{\prime}\) at fixed pressure.  The density-response
terms cancel in the latter derivative by
\eqref{eq:intrinsic-pressure-free-energy}.  The double-prime stress differentiates \(W^{\prime\prime}\) at
fixed pressure \citep[eqs.~(34)--(36)]{fosterxu2025}:
\begin{align}
 \mbf P^{\prime}
 &=\left.\pder{W_s}{\mbf F}\right|_{\bar J}
 ,\label{eq:single-prime-fixed-density-conjugate}\\
 \mbf P^{\prime}&=\left.\pder{W^{\prime}}{\mbf F}\right|_p,
 \label{eq:single-prime-energy-conjugate}\\
 \mbf P^{\prime\prime}
 &=\left.\pder{W^{\prime\prime}}{\mbf F}\right|_p.
 \label{eq:double-prime-energy-conjugate}
\end{align}
Internal variables remain fixed in these constitutive derivatives.

\subsection{The finite-deformation Biot coefficient}
\label{sec:biot-coefficient-and-momentum}

For the isotropic volume response considered here, the fixed-state dependence
of \(\bar J\) on deformation is through \(J\).  Differentiating
\eqref{eq:pressure-legendre-potential} with respect to $\mbf{F}$ gives
\begin{equation}
 \mbf P^{\prime}
 =\mbf P^{\prime\prime}
 +\phi_{s0}p\left.\pder{\bar J}{J}\right|_p
 J\mbf F^{-T}.
 \label{eq:single-prime-first-piola-mineral-transform}
\end{equation}
The corresponding Biot coefficient is
\begin{equation}
 B=1-\phi_{s0}\left.\pder{\bar J}{J}\right|_p.
 \label{eq:finite-deformation-biot-coefficient}
\end{equation}
Our earlier work expresses the same coefficient in terms of the intrinsic
solid specific volume \(\bar v_s=1/\bar\rho_s\) as
\(B=1-(1/v_0^s)\left.\partial\bar v_s/\partial J\right|_p\), where
\(v_0^s=1/\rho_{s0}\) and its solid Jacobian is the present \(J\)
\citep[eq.~(39)]{fosterxu2025}.  Because
\(\bar J=\bar\rho_{s0}\bar v_s\) and
\(\rho_{s0}=\phi_{s0}\bar\rho_{s0}\), these identities give
\((1/v_0^s)\partial\bar v_s/\partial J
=\phi_{s0}\partial\bar J/\partial J\).  In these derivatives, pore pressure
and the current internal state are held fixed.  The referential solid mass
\(\rho_{s0}=J\rho_s\) is already a fixed material parameter by solid mass
conservation, so the two forms agree.  The
Biot coefficient measures the fraction of a skeleton-volume variation
available to change pore volume under those restrictions.  An incompressible
mineral has \(\bar J=1\) and \(B=1\).

Pushing the stresses forward to the current configuration gives
\begin{align}
 \bs\sigma^{\prime}
 &=\bs\sigma^{\prime\prime}+(1-B)p\mbf I,
 \label{eq:single-double-prime-stress-relation}\\
 \bs\sigma
 &=\bs\sigma^{\prime}-p\mbf I,
 \label{eq:total-cauchy-single-prime}\\
 \bs\sigma
 &=\bs\sigma^{\prime\prime}-Bp\mbf I.
 \label{eq:total-cauchy-stress}
\end{align}
Here \(\bs\sigma^{\prime}=J^{-1}\mbf P^{\prime}\mbf F^T\), with the same
push-forward for the double-prime stress, and \(\mbf I\) is the identity
tensor.  Substituting the total first Piola stress directly into the
quasi-static momentum balance gives
\begin{equation}
 \nabla_{\mbf X}\mathbin{\cdot}
 \left(\mbf P^{\prime\prime}-BpJ\mbf F^{-T}\right)
 +J\rho\mbf g=\mbf0,
 \label{eq:reference-momentum-balance}
\end{equation}
where \(\rho\) is total mixture density and \(\mbf g\) is body-force
acceleration.  The Biot coefficient therefore determines both the pressure
term in total stress and the single-prime stress used below to drive plastic
flow.

\section{Poroplastic constitutive response}
\label{sec:inelastic-plastic-driving-stress}

We choose an isotropic model to examine how permanent changes in pore volume
affect pressure coupling.  The logarithmic skeleton and mineral laws give a
closed-form Biot coefficient at each plastic state.  The mineral equation of
state (EOS), yield surface, and dilation rule below specify the material
response.

\subsection{Plastic distention and mineral compression}

We specialize the solid mixture energy to one nonreacting constituent at
constant temperature, with spherical distention and volume-preserving
(isochoric) true-plastic flow.  The elastic and plastic factors satisfy
\begin{equation}
 \bar{\mbf F}=\bar{\mbf F}^e\bar{\mbf F}^p,
 \qquad \mbf A=\mbf A^e\mbf A^p,
 \qquad \det\bar{\mbf F}^p=1,
 \qquad
 \mbf A^e=(a^e)^{1/3}\mbf I,
 \qquad \mbf A^p=(a^p)^{1/3}\mbf I.
 \label{eq:inelastic-kinematic-splits}
\end{equation}
Because the plastic distention factor is spherical, the factors combine to
give
\begin{equation}
 \mbf F=\mbf F^e\mbf F^p,
 \qquad \mbf F^e=(a^e)^{1/3}\bar{\mbf F}^e,
 \qquad \mbf F^p=(a^p)^{1/3}\bar{\mbf F}^p,
 \qquad \det\mbf F^p=a^p.
 \label{eq:aggregate-plastic-split}
\end{equation}
Consequently, \(J^e=\det\mbf F^e=J/a^p\).  Mass conservation of the solid
material and the volume-preserving true-plastic flow give
\(\det\bar{\mbf F}^e=\bar J=\bar\rho_{s0}/\bar\rho_s\), so that
\(J^e=a^e\bar J\).  Elastic skeleton volume is thus divided between
distention of the porous structure and compression of the mineral.

\subsection{Energy-consistent skeleton and mineral response}
\label{sec:elastic-verification-closure}
\label{sec:numerical-solid-eos}

The constitutive construction must reproduce the drained skeleton response
and the mineral compression law through the same stored energy.  Let \(K\)
and \(K_s\) be their reference bulk moduli, with
\(0<K<\phi_{s0}K_s\).  We prescribe the drained volumetric energy per unit
mixture reference volume as \(K(\ln J^e)^2/2\), and the mineral
energy per unit reference mineral volume as
\(K_s(\ln\bar J)^2/2\).  We first identify the stresses conjugate to these
energies, then match their contributions and impose pressure equilibrium.

Here \(\delta\) denotes an infinitesimal virtual variation about the current
state.  We evaluate reversible elastic work with the plastic deformation
held fixed, \(\delta\mbf F^p=\mbf0\), so \(\delta a^p=0\).
The current elastic deformation may be finite; only the variation is
infinitesimal.

The tensor \(\delta\mbf F^e(\mbf F^e)^{-1}\) describes the virtual
deformation relative to the current elastic configuration.  Its symmetric
part is the virtual strain, and its skew part is the virtual rotation.  The
Jacobian variation gives
\begin{align}
 \delta\ln J^e&=\frac{\delta J^e}{J^e},
 \label{eq:elastic-logarithmic-volume-variation}\\
 &=\tr\left[\delta\mbf F^e(\mbf F^e)^{-1}\right].
 \label{eq:elastic-volume-variation-trace}
\end{align}
A purely volumetric virtual variation has equal normal strains in all three
directions, zero shear strain, and zero rotation.  Its virtual
deformation tensor is therefore spherical.  Since \(\tr\mbf I=3\),
\begin{equation}
 \delta\mbf F^e(\mbf F^e)^{-1}
 =\frac13\delta\ln J^e\,\mbf I.
 \label{eq:elastic-spherical-deformation-variation}
\end{equation}
The factor \(1/3\) distributes the fractional volume change equally among
the three directions.  With \(J=a^pJ^e\) and fixed plastic deformation,
\begin{equation}
 \delta\ln J^e=\delta\ln J,
 \qquad \delta a^p=0.
 \label{eq:fixed-plastic-log-volume-variation}
\end{equation}
Multiplying \eqref{eq:elastic-spherical-deformation-variation} on the
right by \(\mbf F^e\) gives
\begin{equation}
 \delta\mbf F^e=\frac13\mbf F^e\,\delta\ln J^e.
 \label{eq:elastic-volumetric-deformation-variation}
\end{equation}
Then taking the variation of \(\mbf F=\mbf F^e\mbf F^p\) at fixed
\(\mbf F^p\) gives
\begin{align}
 \delta\mbf F&=\delta\mbf F^e\mbf F^p,
 \label{eq:fixed-plastic-deformation-variation}\\
 &=\frac13\mbf F\,\delta\ln J^e.
 \label{eq:total-volumetric-deformation-variation}
\end{align}
The variation \eqref{eq:elastic-volumetric-deformation-variation}
preserves the normalized elastic shape \((J^e)^{-1/3}\mbf F^e\):
\begin{align}
 \delta\left[(J^e)^{-1/3}\mbf F^e\right]
 &=(J^e)^{-1/3}\left(\delta\mbf F^e
 -\frac13\mbf F^e\,\delta\ln J^e\right),
 \label{eq:elastic-shape-variation-expansion}\\
 &=\mbf0.
 \label{eq:elastic-shape-variation-zero}
\end{align}
Thus, the elastic volume can vary, \(\delta J^e\ne0\), while the normalized
elastic shape remains fixed.

Hold pore pressure constant during the volume variation, so \(\delta p=0\).
The prescribed pressure may be nonzero.
The mineral volume \(\bar J\) adjusts to satisfy local pressure equilibrium
\eqref{eq:intrinsic-pressure-free-energy}
as \(J^e\) changes.  Both \(\bar J\) and \(a^e=J^e/\bar J\) can therefore
vary during this operation.  At fixed pressure and plastic state, the
energy variation is
\(\delta W^{\prime\prime}=\mbf P^{\prime\prime}:\delta\mbf F\), by
\eqref{eq:double-prime-energy-conjugate}.  Substituting
\eqref{eq:total-volumetric-deformation-variation} gives
\begin{align}
 \delta W^{\prime\prime}
 &=\mbf P^{\prime\prime}:\delta\mbf F,
 \label{eq:drained-energy-first-piola-work}\\
 &=\frac13\left(\mbf P^{\prime\prime}:\mbf F\right)
 \delta\ln J^e,
 \label{eq:drained-energy-contraction-work}\\
 &=\frac13\tr\left(\mbf P^{\prime\prime}\mbf F^T\right)
 \delta\ln J^e,
 \label{eq:drained-energy-volumetric-work}\\
 &=\frac J3\tr\bs\sigma^{\prime\prime}\,\delta\ln J^e.
 \label{eq:drained-energy-cauchy-work}
\end{align}
The double contraction
\(\mbf P^{\prime\prime}:\mbf F= P^{\prime\prime}_{ij}F_{ij}\)
equals \(\tr(\mbf P^{\prime\prime}\mbf F^T)\).  Equation
\eqref{eq:drained-energy-cauchy-work} uses the Cauchy-stress transformation
\(\bs\sigma^{\prime\prime}=J^{-1}\mbf P^{\prime\prime}\mbf F^T\).
Rearranging this transformation defines the double-prime Kirchhoff stress:
\begin{equation}
 \bs\tau^{\prime\prime}\equiv\mbf P^{\prime\prime}\mbf F^T
 =J\bs\sigma^{\prime\prime}.
 \label{eq:drained-kirchhoff-volume-convention}
\end{equation}
The factor \(J\) converts work per unit current mixture volume to work per
unit original reference volume.  The elastic ratio \(J^e\) measures the
recoverable volume change in the energy.

The drained reference response follows from setting the constant pore
pressure to zero.  This is the reference condition for the prescribed energy
\(K(\ln J^e)^2/2\), and the double-prime stress requires only that pressure
be held fixed during the variation.  At \(p=0\) the isochoric contribution is
constant under this volumetric variation, so direct variation gives
\begin{align}
 \delta W^{\prime\prime}
 &=\delta\left[\frac K2(\ln J^e)^2\right],
 \label{eq:drained-logarithmic-energy-variation}\\
 &=K\ln J^e\,\delta\ln J^e.
 \label{eq:drained-logarithmic-energy-variation-result}
\end{align}
Reversible work must equal the change in stored energy for every admissible
virtual volumetric variation.  Equating the coefficients of \(\delta\ln J^e\)
in \eqref{eq:drained-energy-cauchy-work} and
\eqref{eq:drained-logarithmic-energy-variation-result} yields
\begin{equation}
 \left.\frac13\tr\bs\tau^{\prime\prime}\right|_{p=0}
 =K\ln J^e.
 \label{eq:drained-volumetric-energy-derivative}
\end{equation}
Thus the drained mean skeleton Cauchy stress is \(K\ln J^e/J\).
In the elastic specialization, \(a^p=1\) and \(J^e=J\), this becomes
\(K\ln J/J\).

For the mineral, we prescribe the volumetric energy per unit reference
mineral volume as
\begin{equation}
 \bar W_s(\bar J)=\frac{K_s}{2}(\ln\bar J)^2.
 \label{eq:mineral-volumetric-reference-energy}
\end{equation}
To identify its conjugate stress, consider a purely volumetric virtual variation
of the true elastic deformation, holding the true-plastic factor and
elastic distention \(a^e\) fixed.  Applying the determinant identity
\eqref{eq:elastic-volume-variation-trace} and spherical-variation condition
\eqref{eq:elastic-spherical-deformation-variation} to the true elastic
deformation gives
\begin{equation}
 \delta\bar{\mbf F}^e(\bar{\mbf F}^e)^{-1}
 =\frac13\frac{\delta\bar J}{\bar J}\mbf I.
 \label{eq:mineral-spherical-deformation-variation}
\end{equation}
Multiplying \eqref{eq:mineral-spherical-deformation-variation} on the
right by \(\bar{\mbf F}^e\) gives
\begin{equation}
 \delta\bar{\mbf F}^e
 =\frac13\bar{\mbf F}^e\frac{\delta\bar J}{\bar J}.
 \label{eq:mineral-volumetric-deformation-variation}
\end{equation}
The variation \eqref{eq:mineral-volumetric-deformation-variation} changes
the mineral volume while preserving its normalized elastic shape,
\(\bar J^{-1/3}\bar{\mbf F}^e\), by the same cancellation shown in
\eqref{eq:elastic-shape-variation-expansion}.
This variation changes \(J^e=a^e\bar J\) together with \(\bar J\), and
pressure is allowed to change.  The drained skeleton variation instead
fixes pressure and allows elastic distention to adjust.
Let \(\bar{\bs\sigma}_s\) denote the intrinsic solid Cauchy stress.
Its contraction with
\(\delta\bar{\mbf F}^e(\bar{\mbf F}^e)^{-1}\) is work per unit current
mineral volume.  Multiplication by \(\bar J\) converts that work to the
reference mineral volume used by \(\bar W_s\).
Equating reversible work and stored-energy change gives
\begin{align}
 \delta\bar W_s
 &=\bar J\bar{\bs\sigma}_s:
 \left[\delta\bar{\mbf F}^e(\bar{\mbf F}^e)^{-1}\right],
 \label{eq:mineral-volumetric-stress-work}\\
 &=\bar J\bar{\bs\sigma}_s:
 \left(\frac13\mbf I\frac{\delta\bar J}{\bar J}\right),
 \label{eq:mineral-volumetric-work-substitution}\\
 &=\frac13\tr\bar{\bs\sigma}_s\,\delta\bar J.
 \label{eq:mineral-volumetric-work-trace}
\end{align}
Direct variation of \eqref{eq:mineral-volumetric-reference-energy} gives
\begin{align}
 \delta\bar W_s
 &=\frac{K_s}{2}\,2\ln\bar J\frac{\delta\bar J}{\bar J},
 \label{eq:mineral-logarithmic-energy-variation}\\
 &=\frac{K_s\ln\bar J}{\bar J}\,\delta\bar J.
 \label{eq:logarithmic-mineral-energy-derivative}
\end{align}
Equating the coefficients of \(\delta\bar J\) in
\eqref{eq:mineral-volumetric-work-trace} and
\eqref{eq:logarithmic-mineral-energy-derivative} yields
\begin{equation}
 \frac13\tr\bar{\bs\sigma}_s=\frac{K_s\ln\bar J}{\bar J}.
 \label{eq:matched-logarithmic-mineral-stress}
\end{equation}
The mean intrinsic stress is negative in compression, where \(\bar J<1\).
Equation \eqref{eq:matched-logarithmic-mineral-stress} prescribes the
intrinsic mean-stress response.

The tangent of \eqref{eq:matched-logarithmic-mineral-stress} with respect to
\(\ln\bar J\) equals \(K_s\) at \(\bar J=1\), so \(K_s\) is the reference
bulk modulus.  The mean intrinsic stress describes mineral compression, while
the pressure \(\bar p_s=p\) is defined by the energy derivative in
\eqref{eq:local-pressure-equilibrium}.

The stress laws \eqref{eq:drained-volumetric-energy-derivative} and
\eqref{eq:matched-logarithmic-mineral-stress} are connected by the solid's contribution to mixture
stress.  The fluid contributes the intrinsic stress \(-p\mbf I\) and occupies
volume fraction \(1-\phi_s\).  Adding the volume-weighted phase stresses
gives
\begin{equation}
 \bs\sigma=\phi_s\bar{\bs\sigma}_s-(1-\phi_s)p\mbf I.
 \label{eq:constitutive-phase-stress-sum}
\end{equation}
Substituting \eqref{eq:constitutive-phase-stress-sum} into
\eqref{eq:total-cauchy-single-prime} gives
\begin{align}
 \bs\sigma^{\prime}
 &=\phi_s\bar{\bs\sigma}_s-(1-\phi_s)p\mbf I+p\mbf I,
 \label{eq:constitutive-single-prime-phase-substitution}\\
 &=\phi_s\left(\bar{\bs\sigma}_s+p\mbf I\right).
 \label{eq:constitutive-single-prime-phase-stress}
\end{align}
Taking one third of the trace of
\eqref{eq:constitutive-single-prime-phase-stress} and using \(\tr\mbf I=3\) yields
\begin{equation}
 \frac13\tr\bs\sigma^{\prime}
 =\phi_s\left(\frac13\tr\bar{\bs\sigma}_s+p\right).
 \label{eq:constitutive-single-prime-phase-trace}
\end{equation}
Multiplying \eqref{eq:constitutive-single-prime-phase-trace} by \(J\), substituting
\eqref{eq:matched-logarithmic-mineral-stress}, and using solid mass
conservation \eqref{eq:distension}, \(J\phi_s=\phi_{s0}\bar J\), gives
\begin{align}
 \frac13\tr\bs\tau^{\prime}
 &=J\phi_s\left(\frac13\tr\bar{\bs\sigma}_s+p\right),
 \label{eq:constitutive-kirchhoff-intrinsic-stress-trace}\\
 &=J\phi_s\left(p+\frac{K_s\ln\bar J}{\bar J}\right),
 \label{eq:constitutive-kirchhoff-mineral-substitution}\\
 &=\phi_{s0}\bar J
 \left(p+\frac{K_s\ln\bar J}{\bar J}\right),
 \label{eq:constitutive-kirchhoff-solid-mass-substitution}\\
 &=\phi_{s0}\left(p\bar J+K_s\ln\bar J\right).
 \label{eq:constitutive-single-prime-mineral-trace}
\end{align}
At fixed plastic state and elastic shape, the fixed-density stress derivative
\eqref{eq:single-prime-fixed-density-conjugate} and pressure relation
\eqref{eq:intrinsic-pressure-free-energy} give the energy variation in the
independent arguments \((J^e,\bar J)\):
\begin{align}
 \delta W_s
 &=\frac13\tr\bs\tau^{\prime}\,\delta\ln J^e
 -\phi_{s0}p\,\delta\bar J,
 \label{eq:coupled-volume-energy-work}\\
 &=\frac13\tr\bs\tau^{\prime}\,\delta\ln a^e
 +\phi_{s0}\frac13\tr\bar{\bs\sigma}_s\,\delta\bar J.
 \label{eq:distention-mineral-energy-work}
\end{align}
Using \(\delta\ln J^e=\delta\ln a^e+\delta\bar J/\bar J\) in the first
line and substituting \eqref{eq:matched-logarithmic-mineral-stress} and
\eqref{eq:constitutive-single-prime-mineral-trace} gives the second line.
The pressure contributions cancel, so \(\tr\bs\tau^{\prime}/3\) is work
conjugate to \(\ln a^e\), and \(\phi_{s0}\tr\bar{\bs\sigma}_s/3\) is work
conjugate to \(\bar J\).
At fixed distention, mineral work per reference mixture volume is therefore
\(\phi_{s0}\) times mineral work per reference mineral volume.  This is
the energy equivalence used in \eqref{eq:mineral-volumetric-stress-work}.

Because the intrinsic mean stress in
\eqref{eq:matched-logarithmic-mineral-stress} depends only on \(\bar J\),
integration of \eqref{eq:distention-mineral-energy-work} at fixed \(a^e\)
gives the mineral term \(\phi_{s0}K_s(\ln\bar J)^2/2\).  The remaining
volumetric energy depends on \(a^e\) alone.  Its derivative with respect
to \(\ln a^e\) must equal \(\tr\bs\tau^{\prime}/3\).
At \(p=0\), the two effective stresses coincide.  Equating
\eqref{eq:drained-volumetric-energy-derivative} and
\eqref{eq:constitutive-single-prime-mineral-trace} yields
\begin{align}
 K\ln J^e&=\phi_{s0}K_s\ln\bar J,
 \label{eq:drained-mineral-volume-matching}\\
 \ln a^e&=\left(1-\frac{K}{\phi_{s0}K_s}\right)\ln J^e.
 \label{eq:drained-distention-volume-matching}
\end{align}
Using \eqref{eq:drained-distention-volume-matching} to express the drained
mean Kirchhoff stress \eqref{eq:drained-volumetric-energy-derivative} in
terms of \(\ln a^e\) and integrating fixes the distention energy.  With zero energy at
\(a^e=\bar J=1\), the combined volumetric energy is
\begin{equation}
 W_{s,\mathrm{vol}}(J^e,\bar J)
 =\frac{K}{2\left(1-K/(\phi_{s0}K_s)\right)}
 \left[\ln\left(\frac{J^e}{\bar J}\right)\right]^2
 +\frac{\phi_{s0}K_s}{2}(\ln\bar J)^2.
 \label{eq:distention-mineral-volumetric-energy}
\end{equation}
Here \(a^e=J^e/\bar J\) has been substituted so that the energy is expressed
in the independent arguments \((J^e,\bar J)\).  At equilibrium under zero pressure, substitution of
\eqref{eq:drained-mineral-volume-matching}--\eqref{eq:drained-distention-volume-matching}
into \eqref{eq:distention-mineral-volumetric-energy} recovers the prescribed
drained skeleton energy \(K(\ln J^e)^2/2\).

The energy \eqref{eq:distention-mineral-volumetric-energy} determines mineral volume at a prescribed pore
pressure through \eqref{eq:intrinsic-pressure-free-energy}.  To evaluate
that derivative, vary \(\bar J\) at fixed \(J^e\), elastic shape,
and plastic state.  These constraints fix \(\mbf F^e\) and hence
\(\mbf F\), while \(a^e=J^e/\bar J\) changes.  Both volumetric energy
terms therefore contribute.  Differentiating
\eqref{eq:distention-mineral-volumetric-energy} gives
\begin{equation}
 \left.\pder{W_{s,\mathrm{vol}}}{\bar J}\right|_{J^e}
 =-\frac{K\ln(J^e/\bar J)}
 {\left(1-K/(\phi_{s0}K_s)\right)\bar J}
 +\frac{\phi_{s0}K_s\ln\bar J}{\bar J}.
 \label{eq:matched-energy-density-derivative}
\end{equation}
Applying the pressure-equilibrium relation
\eqref{eq:intrinsic-pressure-free-energy} sets the derivative in
\eqref{eq:matched-energy-density-derivative} equal to
\(-\phi_{s0}p\).  Rearranging gives
\begin{equation}
 p\bar J+K_s\ln\bar J
 =\frac{KK_s}{\phi_{s0}K_s-K}\ln\left(\frac{J^e}{\bar J}\right).
 \label{eq:constitutive-characteristic-volume-relation}
\end{equation}
Collecting the logarithmic terms in
\eqref{eq:constitutive-characteristic-volume-relation} yields the implicit
mineral EOS residual
\begin{equation}
 \boxed{
 K_s\ln\bar J
 +\left(1-\frac{K}{\phi_{s0}K_s}\right)p\bar J
 -\frac{K}{\phi_{s0}}\ln J^e=0.}
 \label{eq:verification-mineral-factors}
\end{equation}
For virgin elasticity, \(J^e=J\), this residual also follows by eliminating
the mineral deformation factors from Gajo's logarithmic constitutive
relations \citep[Eqs.~(3.27), (3.32), and (3.34)]{gajo2010}, with his
\(J_s=\bar J\) and \(1-n_0=\phi_{s0}\).
Once the positive root of \eqref{eq:verification-mineral-factors} is found,
mass conservation \eqref{eq:distension} gives
\(\bar\rho_s=\bar\rho_{s0}/\bar J\) and
\(\phi_s=\phi_{s0}\bar J/J\).  Written as local pressure equilibrium,
\eqref{eq:verification-mineral-factors} becomes
\begin{equation}
 \bar p_s(\bar\rho_s,J^e)
 =\frac{(K/\phi_{s0})\ln J^e-K_s\ln\bar J}
 {\left(1-K/(\phi_{s0}K_s)\right)\bar J}=p.
 \label{eq:elastic-mineral-pressure}
\end{equation}
For \(p\geq0\), the residual in
\eqref{eq:verification-mineral-factors} is strictly increasing from
negative to positive infinity as \(\bar J\) increases, so its positive
root is unique.  At \(p=0\), \eqref{eq:verification-mineral-factors} gives
\(K_s\ln\bar J=(K/\phi_{s0})\ln J^e\), or
\(\bar J=(J^e)^{K/(\phi_{s0}K_s)}\).  For \(p<0\), multiple positive roots
can exist; at fixed \(J^e\), we select the solution that varies continuously
from this zero-pressure root as pressure decreases.

We retain this solution only while the residual in
\eqref{eq:verification-mineral-factors} has a
positive derivative with respect to \(\bar J\), which requires
\begin{equation}
 K_s+\left(1-\frac{K}{\phi_{s0}K_s}\right)p\bar J>0.
 \label{eq:trace-mineral-stability-domain}
\end{equation}
Condition \eqref{eq:trace-mineral-stability-domain} keeps the scalar
derivative positive.  We also require \(0<\bar J<e\) for positive mineral
stiffness and \(0<\phi_s<1\) for positive solid and pore volumes.


\subsection{Closed-form Biot coefficient from implicit differentiation}

The derivative defining \(B\) in \eqref{eq:finite-deformation-biot-coefficient}
fixes pore pressure and the current plastic configuration; the referential
solid mass is already fixed by mass conservation.  Because \(J^e=J/a^p\),
\(\partial\ln J^e/\partial J=1/J\) under this constraint.
Differentiating \eqref{eq:verification-mineral-factors} therefore gives
\begin{align}
 \left[\frac{K_s}{\bar J}
 +\left(1-\frac{K}{\phi_{s0}K_s}\right)p\right]
 \left.\pder{\bar J}{J}\right|_{p,\mbf F^p}
 &=\frac{K}{\phi_{s0}J},
 \label{eq:solid-pressure-implicit-tangent}\\
 \left.\pder{\bar J}{J}\right|_{p,\mbf F^p}
 &=\frac{K\bar J}
 {\phi_{s0}J\left[K_s+\left(1-\dfrac{K}{\phi_{s0}K_s}\right)p\bar J\right]}.
 \label{eq:solid-density-eos-tangent}
\end{align}
Substituting \eqref{eq:solid-density-eos-tangent} into
\eqref{eq:finite-deformation-biot-coefficient} yields
\begin{equation}
 \boxed{B=1-\frac{K\bar J}
 {J\left[K_s+\left(1-\dfrac{K}{\phi_{s0}K_s}\right)p\bar J\right]}.}
 \label{eq:poroplastic-biot-correction}
\end{equation}
The Biot coefficient can therefore be evaluated directly from \(J\), \(p\),
and the mineral-volume ratio \(\bar J\) obtained by solving
\eqref{eq:verification-mineral-factors}.
Plastic history enters through
\(J^e=J/a^p\) in \eqref{eq:verification-mineral-factors}, whereas the
denominator of \eqref{eq:poroplastic-biot-correction} is expressed in the
total volume, consistent with the derivative in
\eqref{eq:finite-deformation-biot-coefficient}.

Solid mass conservation \eqref{eq:distension} also gives \(\bar J=J\phi_s/\phi_{s0}\).
Substituting this identity into \eqref{eq:poroplastic-biot-correction}
and clearing the nested fractions gives the solid-volume-fraction form
\begin{equation}
 B=1-\frac{KK_s\phi_{s0}\phi_s}
 {\phi_{s0}^2K_s^2+(\phi_{s0}K_s-K)pJ\phi_s}.
 \label{eq:biot-solid-volume-fraction-form}
\end{equation}
Let \(\phi=1-\phi_s\) be the current porosity and
\(\phi_0=1-\phi_{s0}\) its reference value.  The equivalent porosity form is
\begin{equation}
 B=1-\frac{KK_s(1-\phi_0)(1-\phi)}
 {(1-\phi_0)^2K_s^2+\left[(1-\phi_0)K_s-K\right]pJ(1-\phi)}.
 \label{eq:biot-porosity-form}
\end{equation}
Equations \eqref{eq:biot-solid-volume-fraction-form} and
\eqref{eq:biot-porosity-form} use the current phase volumes determined by
\eqref{eq:verification-mineral-factors} and \eqref{eq:distension}, and retain
the fixed-plastic-state definition of \(B\).

For a virgin state, one that has not yielded, \(a^p=1\) and \(J^e=J\).  We
denote the Biot coefficient evaluated with this mineral state by
\(B_{\mathrm{el}}(p,J)\).
At \(J=\bar J=1\) and \(p=0\), \eqref{eq:poroplastic-biot-correction} recovers
\(B_0=1-K/K_s\) \citep{biotwillis1957}.  At fixed total volume and pressure,
increased plastic distention reduces \(J^e\) and \(\bar J\) while
\eqref{eq:trace-mineral-stability-domain} holds, and hence increases \(B\).  This dependence follows from the
mineral equation \eqref{eq:verification-mineral-factors} and the Biot
coefficient \eqref{eq:poroplastic-biot-correction}.

\subsection{Coupled stored energy and stress}

Adding a neo-Hookean isochoric response of shear modulus \(G\) to the
volumetric energy \eqref{eq:distention-mineral-volumetric-energy} gives the
complete solid stored energy per unit original mixture reference volume,
\begin{align}
 W_s(\mbf F^e,\bar J)
 &=\frac{G}{2}\left[(J^e)^{-2/3}
 \tr\left(\mbf F^e(\mbf F^e)^T\right)-3\right] \notag\\
 &\quad+\frac{K}{2\left(1-K/(\phi_{s0}K_s)\right)}
 \left[\ln\left(\frac{J^e}{\bar J}\right)\right]^2
 +\frac{\phi_{s0}K_s}{2}(\ln\bar J)^2.
 \label{eq:coupled-skeleton-mineral-energy}
\end{align}
The first term governs isochoric deformation, while the remaining terms are
the distention and mineral energies derived above.  Differentiation at fixed
\(\bar J\) gives the single-prime Kirchhoff stress.  Using the mineral EOS
\eqref{eq:verification-mineral-factors} in the second line gives
\begin{align}
 \bs\tau^{\prime}
 &=G(J^e)^{-2/3}
 \operatorname{dev}\left(\mbf F^e(\mbf F^e)^T\right)
 +\frac{K}{1-K/(\phi_{s0}K_s)}
 \ln\left(\frac{J^e}{\bar J}\right)\mbf I,
 \label{eq:single-prime-stress-from-energy}\\
 &=G(J^e)^{-2/3}
 \operatorname{dev}\left(\mbf F^e(\mbf F^e)^T\right)
 +K\left(\ln J^e+\frac{p\bar J}{K_s}\right)\mbf I.
 \label{eq:matched-logarithmic-single-prime-stress}
\end{align}
Using the fixed-density/fixed-pressure stress relation
\eqref{eq:single-double-prime-stress-relation} and
\(\mbf P^{\prime\prime}=J\bs\sigma^{\prime\prime}\mbf F^{-T}\) gives the
double-prime first Piola stress
\begin{align}
 \mbf P^{\prime\prime}
 &=G(J^e)^{-2/3}
 \operatorname{dev}\left(\mbf F^e(\mbf F^e)^T\right)\mbf F^{-T} \notag\\
 &\quad+\left\{K\ln J^e
 +\left[\frac{K\bar J}{K_s}-(1-B)J\right]p\right\}\mbf F^{-T}.
 \label{eq:matched-logarithmic-double-prime-stress}
\end{align}
Here \(\bar J\) is the positive root of
\eqref{eq:verification-mineral-factors}.  The total first Piola stress used in
the momentum balance and constitutive computations is
\(\mbf P=\mbf P^{\prime\prime}-BpJ\mbf F^{-T}\), as in
\eqref{eq:reference-momentum-balance}.  Its push-forward gives
\(\bs\sigma=\bs\sigma^{\prime\prime}-Bp\mbf I\), with \(B\) given by
\eqref{eq:poroplastic-biot-correction}.

The term \((K/K_s)p\bar J\) in
\eqref{eq:matched-logarithmic-single-prime-stress} expresses the contribution
of mineral compression to the mean effective stress.  It remains active in
purely elastic deformation, where \(\mbf F^p=\mbf I\).  In that case, the
total Kirchhoff stress agrees with Gajo's logarithmic elastic specialization
\citep[eq.~(3.47)]{gajo2010}.  Gajo's pressure factor in his equation~(3.48)
is defined relative to the drained elastic stress.  Here \(B\) is the
fixed-pressure mineral-volume derivative, and
\eqref{eq:matched-logarithmic-double-prime-stress} retains the corresponding
pressure correction in the double-prime stress.  These definitions distinguish
the coefficient multiplying pressure relative to a drained stress from the
coefficient measuring the incremental pressure response.
For the logarithmic elastic specialization, the coupling factor in Gajo's
rate equations~(5.12) and~(5.14) reduces to the present \(B\).

The small-strain elastic limit recovers the classical Biot effective-stress
relation.  First, pushing forward
\eqref{eq:matched-logarithmic-double-prime-stress} gives
\begin{align}
 \bs\sigma^{\prime\prime}
 &=J^{-1}\mbf P^{\prime\prime}\mbf F^T,
 \label{eq:double-prime-cauchy-push-forward}\\
 &=\frac{G}{J}(J^e)^{-2/3}
 \operatorname{dev}\left(\mbf F^e(\mbf F^e)^T\right) \notag\\
 &\quad+\left\{\frac{K}{J}\ln J^e
 +\left[\frac{K\bar J}{JK_s}-(1-B)\right]p\right\}\mbf I.
 \label{eq:matched-logarithmic-double-prime-cauchy-stress}
\end{align}
Linearize about the undeformed, zero-pressure state with
\(\mbf F^p=\mbf I\), so \(\mbf F^e=\mbf F\), and define the infinitesimal
strain by
\begin{equation}
 \bs\varepsilon=\frac12\left[
 \nabla_{\mbf X}\mbf u_s+(\nabla_{\mbf X}\mbf u_s)^T\right].
 \label{eq:biot-limit-infinitesimal-strain}
\end{equation}
To first order in displacement gradient and pressure,
\begin{align}
 \ln J^e&\simeq\tr\bs\varepsilon,
 \label{eq:biot-limit-logarithmic-volume}\\
 \operatorname{dev}\left(\mbf F^e(\mbf F^e)^T\right)
 &\simeq2\operatorname{dev}\bs\varepsilon.
 \label{eq:biot-limit-isochoric-strain}
\end{align}
The coefficient multiplying \(p\) in
\eqref{eq:matched-logarithmic-double-prime-cauchy-stress} is zero at
\(J=\bar J=1\), since \(1-B_0=K/K_s\).  Its variation multiplied by
\(p\) is second order.  Likewise, variations of the prefactors multiplying
strain contribute only at second order.  The linearized effective and total
stresses are therefore
\begin{align}
 \bs\sigma^{\prime\prime}
 &\simeq2G\operatorname{dev}\bs\varepsilon
 +K\tr\bs\varepsilon\,\mbf I,
 \label{eq:biot-limit-effective-stress}\\
 \bs\sigma
 &\simeq2G\operatorname{dev}\bs\varepsilon
 +K\tr\bs\varepsilon\,\mbf I-B_0p\mbf I.
 \label{eq:biot-limit-total-stress}
\end{align}
Thus, with tensile stress positive, adding \(B_0p\mbf I\) to the total
stress recovers the drained elastic stress in the linearized theory.

Zero-pressure calibration supplies the drained double-prime skeleton stress.
For the assumed elastic law, these tests determine \(G\) and \(K\): at
\(p=0\), the pressure correction in
\eqref{eq:matched-logarithmic-double-prime-stress} vanishes and the stress
measures coincide.  The mineral modulus \(K_s\) and reference solid fraction
\(\phi_{s0}\) supply additional information for the coupled energy.
That energy and local pressure equilibrium extend the double-prime response
to nonzero pressure.  The effective-stress relations then supply the single-prime
stress driving plastic flow and the total stress entering momentum balance.
Thus the double-prime stress at nonzero pressure includes both the calibrated
drained response and the pressure correction derived from the mineral law.

At fixed deformation and plastic state, with mineral volume in local
equilibrium, the Biot coefficient measures the sensitivity of total stress
to pore pressure:
\begin{equation}
 \left.\pder{\bs\sigma}{p}\right|_{\mbf F,\mbf F^p}
 =-B\mbf I.
 \label{eq:total-stress-pressure-tangent}
\end{equation}

For an elastic specimen in the unjacketed test, where the pore fluid and the
external boundary carry the same hydrostatic pressure, \(J=\bar J\),
\(p=-K_s\ln J/J\), and \(\phi_s=\phi_{s0}\).  Substitution gives the total
stress \(-p\mbf I\).  In the incompressible-mineral limit,
\(K_s\to\infty\), the mineral volume remains fixed, \(B\to1\), and
\(\mbf P^{\prime\prime}\) reduces to the drained skeleton response.

\subsection{Plastic driving stress, yield, and flow}
\label{sec:plastic-loading-incremental-rate}

Plastic flow changes \(\mbf F^e\) at fixed total deformation.  The
single-prime Kirchhoff stress, its Mandel pullback, and the deviatoric stress
in the intermediate configuration are
\begin{align}
 \bs\tau^{\prime}
 &=\mbf P^{\prime}\mbf F^T,
 \label{eq:kirchhoff-stress-definition}\\
 \bs\tau^{\prime}&=\mbf P^{\prime\prime}\mbf F^T+(1-B)pJ\mbf I,
 \label{eq:single-prime-kirchhoff-stress}\\
 \mbf M^{\prime}&=(\mbf F^e)^T
 \bs\tau^{\prime}(\mbf F^e)^{-T},
 \label{eq:complete-mandel-driving-stress}\\
 \mbf S&=\operatorname{dev}\mbf M^{\prime}
 =(\mbf F^e)^T
 \operatorname{dev}\bs\tau^{\prime}(\mbf F^e)^{-T}.
 \label{eq:mandel-driving-stress}
\end{align}
At fixed total deformation, differentiating
\(\mbf F=\mbf F^e\mbf F^p\) gives
\(\dot{\mbf F}^e=-\mbf F^e\dot{\mbf F}^p(\mbf F^p)^{-1}\).
The chain rule therefore contributes
\(-\mbf M^{\prime}:\dot{\mbf F}^p(\mbf F^p)^{-1}\) to the stored-energy
rate and the opposite-signed term to the dissipation.  Thus
\(\mbf M^{\prime}\) is work conjugate to the plastic velocity gradient, as in
standard multiplicative plasticity \citep{yamakawaetal2021}.  The single-prime
stress is obtained by differentiating \(W_s\) with respect to deformation at
fixed intrinsic density.  Equivalently, it is the deformation derivative of
the pressure Legendre transform \(W^{\prime}\) at fixed pore pressure, with
the plastic state held fixed in both derivatives
\eqref{eq:single-prime-fixed-density-conjugate}--\eqref{eq:single-prime-energy-conjugate}.

Drumheller's mixture dissipation analysis writes the true-plastic and
plastic-distention powers per unit current volume using a Cauchy-type stress
\citep[eqs.~(108)--(112)]{drumheller2000}.  In the present single-phase
notation, let
\(\mbf L^p=\dot{\mbf F}^p(\mbf F^p)^{-1}\) and
\(\mbf l^p=\mbf F^e\mbf L^p(\mbf F^e)^{-1}\).  Multiplication by \(J\)
converts the spatial plastic power to power per unit original mixture
reference volume.  Using \(\bs\tau^{\prime}=J\bs\sigma^{\prime}\), this gives
\(J\bs\sigma^{\prime}:\mbf l^p
=\mbf M^{\prime}:\mbf L^p\).  Thus Drumheller's Cauchy-stress form and the
Kirchhoff--Mandel form used here represent the same plastic power in different
volume measures.  The resulting deviatoric and volumetric conjugate stresses
are also obtained directly in the single-phase specialization
\citep[Sec.~3, eq.~(8)]{foster2026nonassociated}.  For the isotropic energy,
\(\mbf M^{\prime}\) and \(\mbf S\) are symmetric.  Define the
compression-positive mean and equivalent shear stresses by
\begin{equation}
 p^{\prime}=-\frac{1}{3}\tr\bs\tau^{\prime},
 \qquad q=\sqrt{\frac{3}{2}\mbf S:\mbf S}.
 \label{eq:finite-plastic-stress-invariants}
\end{equation}
The prime distinguishes the plastic mean-stress argument from pore pressure
\(p\).  We choose the Drucker--Prager surface and dilation rule
\begin{align}
 f&=q-Mp^{\prime}-c\leq0,
 \label{eq:dp-yield-observable}\\
 \dot{\mbf F}^p(\mbf F^p)^{-1}
 &=\dot\gamma\left(\frac{3\mbf S}{2q}+\frac{\beta}{3}\mbf I\right),
 \label{eq:dp-flow}\\
 \dot\gamma&\geq0,
 \qquad \dot\gamma f=0.
 \label{eq:plastic-complementarity}
\end{align}
Here \(M\) is the friction slope, \(c\) is cohesion, and \(\beta\) is the
dilation slope, with \(0\leq\beta\leq M\) and \(c\geq0\).  We consider
the smooth part of the yield cone, where \(q>0\).  To permit isotropic hardening, let
\(\gamma\) be the accumulated plastic multiplier and prescribe
\begin{align}
 c(\gamma)&=c_0+H\gamma,
 \qquad c_0\geq0,\quad H\geq0,
 \label{eq:plastic-isotropic-hardening}\\
 \gamma(t)&=\int_0^t\dot\gamma(s)\dd s,
 \qquad \gamma(0)=0.
 \label{eq:plastic-accumulated-multiplier}
\end{align}
Here \(H\) has units of stress.  This dissipative hardening law increases
the yield strength with plastic history while retaining the nonassociated
flow direction and the stated stored energy.  The choice \(H=0\) recovers
perfect plasticity.  Taking the trace of the flow rule \eqref{eq:dp-flow} gives
\begin{equation}
 \frac{\dot a^p}{a^p}=\beta\dot\gamma.
 \label{eq:reconstruct-pore-allocation}
\end{equation}
The remaining plastic factor is isochoric.  Because
\(\mbf M^{\prime}=\mbf S-p^{\prime}\mbf I\), the plastic dissipation per unit
reference volume is
\begin{align}
 \mc D^p
 &=\mbf M^{\prime}:\left[\dot{\mbf F}^p(\mbf F^p)^{-1}\right] \notag\\
 &=\dot\gamma\left(
 \frac{3}{2q}\mbf S:\mbf S-\beta p^{\prime}\right) \notag\\
 &=\dot\gamma(q-\beta p^{\prime}),
 \label{eq:plastic-power}\\
 \mc D^p&=\dot\gamma\left[c+(M-\beta)p^{\prime}\right]\geq0.
 \label{eq:plastic-dissipation-on-yield}
\end{align}
The first result uses \(\mbf S:\mbf S=2q^2/3\).  On the yield surface, the
inequality follows from the yield condition and the stated parameter
restrictions, with \(c\) evaluated at the current accumulated multiplier.
These equations specify a dilative model with optional isotropic hardening.
Representing pore collapse, a cap, or the cone apex would require additional
constitutive equations for those cases.

For each load increment, subscripts \(n\) and \(n+1\) denote the known
state at the beginning of the increment and the state to be determined at
its end, respectively.  Introduce the symmetric logarithmic plastic increment
\(\bs\Delta^p\) and solve
\begin{align}
 \mbf F^p_{n+1}&=\exp\left(\bs\Delta^p\right)\mbf F^p_n,
 \label{eq:plastic-exponential-update}\\
 \bs\Delta^p&=\Delta\gamma\left(
 \frac{3\mbf S_{n+1}}{2q_{n+1}}+\frac{\beta}{3}\mbf I\right),
 \label{eq:implicit-plastic-flow-increment}\\
 f_{n+1}&=0
 \quad\text{during plastic flow}.
 \label{eq:discrete-consistency}
\end{align}
This exponential update evaluates the flow direction at the returned state.
The accumulated multiplier is updated by
\begin{equation}
 \gamma_{n+1}=\gamma_n+\Delta\gamma.
 \label{eq:discrete-plastic-hardening}
\end{equation}
Substitution into \eqref{eq:plastic-isotropic-hardening} places the hardening
contribution inside the same local consistency solve.  It introduces no
additional independent local unknown.
The elastic trial at the current increment holds the plastic history fixed at
\(\mbf F^p_n\) and uses
\(\mbf F^{e,\mathrm{tr}}_{n+1}=\mbf F_{n+1}(\mbf F^p_n)^{-1}\).
The trial stress and \(f^{\mathrm{tr}}_{n+1}\) are evaluated from this
candidate state at \(n+1\).  A trial state inside the yield surface leaves
the plastic history unchanged.  An active trial requires a coupled
solve for the six independent components of \(\bs\Delta^p\) and \(\Delta\gamma\).  At every local iteration,
\(\mbf F^e\), intrinsic density, \(B\), and stress are recomputed from the
same candidate state.  The determinant identity
\(a^p_{n+1}=a^p_n\exp(\beta\Delta\gamma)\) provides an independent check
of this update.
